\documentclass[12pt,a4paper]{article}
\usepackage[utf8]{inputenc}
\usepackage[spanish,es-noquoting]{babel}
\usepackage[T1]{fontenc}
\usepackage{lmodern}
\usepackage{geometry}
\geometry{top=2cm, bottom=2cm, left=2.5cm, right=2.5cm}
\usepackage{xcolor}
\usepackage{listings}
\usepackage{enumitem}
\usepackage{fancyhdr}
\usepackage{hyperref}
\usepackage{tikz}
\usetikzlibrary{arrows.meta, positioning, calc}

\definecolor{codegreen}{rgb}{0,0.6,0}
\definecolor{backcolour}{rgb}{0.95,0.95,0.92}

\lstset{
    language=C,
    backgroundcolor=\color{backcolour},
    commentstyle=\color{codegreen},
    keywordstyle=\color{blue},
    basicstyle=\ttfamily\small,
    breaklines=true,
    frame=single,
    numbers=left,
    tabsize=4,
    extendedchars=true,
    showstringspaces=false,
    literate={á}{{\'a}}1 {é}{{\'e}}1 {í}{{\'i}}1 {ó}{{\'o}}1 {ú}{{\'u}}1 {ñ}{{\~n}}1
}

\pagestyle{fancy}
\fancyhf{}
\lhead{Monitorías Sistemas Operativos 2026-I}
\rhead{Capítulo 0: Matriz Dinámica}
\cfoot{\thepage}

\title{\textbf{Guía de Aprendizaje: Creación de una Matriz Dinámica}}
\author{Malak Sanchez \\ \small Monitorías de Sistemas Operativos}
\date{}

\begin{document}

\maketitle

\section{El reto: ¿Qué es una matriz dinámica?}
En C, una matriz en el Stack se declara como \texttt{int mat[F][C]}. Pero si queremos que el tamaño sea dinámico, no podemos hacerlo de forma tan directa. Necesitamos usar un \textbf{puntero a punteros} (\texttt{int **}).

\section{Estructura Mental}
Imagina una matriz dinámica no como un bloque sólido, sino como:
\begin{enumerate}
    \item Un puntero principal (\texttt{int **matriz}).
    \item Un arreglo de \textbf{punteros} (uno por cada fila).
    \item Cada uno de esos punteros apunta a un arreglo de \textbf{enteros} (las columnas).
\end{enumerate}

\section{Dibujo Ilustrativo}

\begin{center}
\begin{tikzpicture}[node distance=1.5cm, font=\small, >=Stealth]
    \node[draw, fill=blue!10, minimum width=2.5cm, minimum height=1cm] (stack) at (0,3) {\texttt{int **matrix}};
    \node[above=0.1cm of stack] {\textbf{STACK}};

    \node[above=0.3cm] at (4,4.1) {\textbf{HEAP (Level 1)}};
    \foreach \i in {0,1,2} {
        \node[draw, fill=green!10, minimum width=2.8cm, inner sep=4pt] (rowptr\i) at (4, 3.8 - \i*0.7) {\texttt{matrix[\i]}};
    }
    \draw[very thick] (2.5, 4.1) rectangle (5.5, 1.7);

    \node[above=0.3cm] at (9.1,4.1) {\textbf{HEAP (Level 2)}};
    \foreach \r in {0,1,2} {
        \foreach \c in {0,1,2} {
            \node[draw, fill=yellow!10, minimum width=1cm] (data\r\c) at (8+\c*1.1, 3.8 - \r*0.7) {\small Data};
        }
    }

    \draw[->, thick, blue] (stack.east) -- (rowptr0.west);
    \draw[->, thick, red] (rowptr0.east) -- (data00.west);
    \draw[->, thick, red] (rowptr1.east) -- (data10.west);
    \draw[->, thick, red] (rowptr2.east) -- (data20.west);

\end{tikzpicture}
\end{center}
\newpage

\section{Cómo se construye (El Código)}
Debemos hacer la reserva en \textbf{capas}. Primero las filas, luego cada fila individual.

\begin{lstlisting}
#include <stdio.h>
#include <stdlib.h>

int main(){
    int rows = 3, cols = 4;

    // 1. Allocate the array of pointers (the rows)
    int **matrix = (int **) malloc(rows * sizeof(int *));
    if(!matrix){
        printf("Error: Could not allocate memory.\n");
        return 1;
    }

    // 2. For each row, allocate the array of data
    for(int i=0; i < rows; i++){
        matrix[i] = (int *) malloc(cols * sizeof(int));
        if(!matrix[i]){
            printf("Error: Could not allocate memory.\n");
            return 1;
        }
    }

    // 3. Normal usage: matrix[i][j]
    matrix[1][2] = 50;
    
    // 4. FREE (In inverse order of creation)
    for(int i=0; i < rows; i++){
        free(matrix[i]); // First we free each row
    }
    free(matrix); // At the end the "parent" of all
    
    return 0;
}
\end{lstlisting}

\section{Regla de Oro}
Por cada \texttt{malloc} que hagas, debe existir un \texttt{free}.
En una matriz de $N$ filas, haces $N+1$ allocations (\texttt{malloc}): una para el arreglo de punteros y $N$ para las filas datos. Por lo tanto, ¡necesitas $N+1$ llamadas a \texttt{free}!

\section{Errores comunes}
\begin{itemize}
    \item \textbf{Liberar solo el padre:} Si haces \texttt{free(matriz)} sin liberar las filas primero, habrás perdido la dirección de las filas y esa memoria quedará "bloqueada" (Memory Leak).
    \item \textbf{Confusión en sizeof:} En el primer \texttt{malloc} se usa \texttt{sizeof(int *)}, porque estamos guardando \textbf{punteros}. En el segundo se usa \texttt{sizeof(int)}, porque guardamos \textbf{enteros}.
\end{itemize}

\end{document}
